\section{Introduction}
Analysts and investors routinely ask questions such as ``What was Home Depot's revenue in the third quarter of 2022?'' or ``Which of these companies had the highest operating margin in 2021?''. Language models answer such questions poorly from memory, and retrieval-augmented generation (RAG) \cite{lewis2020rag} is the usual remedy: retrieve passages first, then answer from them. Earnings calls make the retrieval step unusually hard. Every quarter repeats the same phrasing with different numbers, so a passage from the wrong quarter looks almost identical to the right one, and an answer built on it is fluent and wrong.

Graph-based RAG organises a corpus as a knowledge graph and retrieves connected facts \cite{edge2024graphrag,guo2024lightrag,gutierrez2024hipporag}. TG-RAG \cite{han2025tgrag} adds time: each relation carries a timestamp, timestamps form a year--quarter--month--day hierarchy, and only relations dated inside the period named in the question may start the retrieval. On its earnings-call benchmark, ECT-QA, this filter is one of the method's main strengths.

The filter is all-or-nothing, however, and it filters on timestamps that a language model wrote while building the graph. In the graph we built, 832 of 8,408 stored facts (9.9\%) carry a timestamp that does not even follow a valid calendar form (``2023 QQ3'', ``Q1 2023'', ``LAST YEAR''). Correct timestamps can also be the wrong granularity: a full-year fact is excluded from a question about one of that year's quarters, and a neighbouring-quarter fact is excluded even when nothing from the exact quarter exists. In one of our development questions, none of the 60 most similar candidates carried the asked quarter, so the filter kept nothing and the model invented three figures.

This paper asks whether softening the filter helps. Our contributions are:
\begin{itemize}
  \item \textbf{A soft temporal kernel} $\kappa(\tau) = \max_{t \in T_q} \exp(-\lambda\, d_H(\tau, t))$ over a calendar tree, with one parameter, which equals the hard filter as $\lambda \to \infty$ (confirmed empirically).
  \item \textbf{Kernel-weighted retrieval}: the kernel is used in seed selection, in the personalized PageRank walk and in evidence scoring, and compared with a question-independent decay in the style of DyG-RAG \cite{sun2025dygrag} and with time-partitioned hybrid BM25 and dense seeding.
  \item \textbf{A timestamp-noise protocol} that shifts a controlled share of stored timestamps deterministically, so that every method sees identical noise.
  \item \textbf{A pre-registered evaluation}: the baseline is re-implemented and verified against the released code, all settings are fixed on a development set, and the test comparisons are written down before the single test run.
\end{itemize}
The main result is negative. On 152 held-out questions the soft kernel ties with the reproduced baseline (0.603 against 0.599 Correct). A per-question analysis shows why: the kernel changes the evidence for almost every question, helping multi-period questions and hurting single-period, multi-company ones. When stored timestamps are noisy, the hard filter degrades and the soft kernel does not; the gap at 30\% noise is significant on its own but not after correcting for multiple tests.
